\documentclass[12pt,a4paper]{article}
\usepackage[utf8]{inputenc}
\usepackage[T5]{fontenc}
\usepackage[vietnamese]{babel}
\usepackage[margin=1.3cm,top=1.0cm,bottom=1.05cm,includehead,headheight=18pt,headsep=8pt]{geometry}
\usepackage{amsmath,amssymb,array,tabularx,enumitem,fancyhdr,tikz,xcolor,needspace}
\usetikzlibrary{calc,angles,quotes,arrows.meta,babel,patterns}
\definecolor{ink}{RGB}{30,55,70}

\newcommand{\headerQuote}{%
  \ifcase\value{page}\relax
    \or Math gives us hope that every problem has a solution% Trang 1
    \or Math + You = Champion% Trang 2
    \or Where Frustration Becomes Inspiration% Trang 3
    \or Turning Confusion into Clarity% Trang 4
    \or Don't be an X, be a Y of mathematics% Trang 5
    \or Don't get mad. Learn to add% Trang 6
    \or The only way to learn mathematics is to do mathematics% Trang 7
    \or Small steps every day lead to big results% Trang 8
    \or Every expert was once a beginner% Trang 9
    \or Success is the sum of small efforts repeated daily% Trang 10
    \or Mathematics is the music of reason% Trang 11
    \or Believe you can and you're halfway there% Trang 12
    \else Math gives us hope that every problem has a solution%
  \fi
}
\pagestyle{fancy}\fancyhf{}
\lhead{\small\textit{\headerQuote}}
\rhead{\small\textit{Người biên soạn: Dương Gia Bảo}}
\rfoot{\small Trang \thepage}
\renewcommand{\headrulewidth}{0.4pt}
\setlength{\headheight}{18pt}
\setlength{\parindent}{0pt}
\setlength{\parskip}{0pt}
\setlength{\tabcolsep}{1.5pt}
\renewcommand{\arraystretch}{1.0}
\newcommand{\parttitle}[1]{\par\vspace{1.5pt}\needspace{2\baselineskip}{\large\bfseries\color{ink}#1}\par\nopagebreak\vspace{1pt}}
\newcommand{\four}[4]{\par\noindent\begin{tabularx}{\linewidth}{@{}*{4}{>{\raggedright\arraybackslash}X}@{}}\textbf{A.} #1&\textbf{B.} #2&\textbf{C.} #3&\textbf{D.} #4\end{tabularx}}
\newcommand{\two}[4]{\par\noindent\begin{tabularx}{\linewidth}{@{}*{2}{>{\raggedright\arraybackslash}X}@{}}\textbf{A.} #1&\textbf{B.} #2\\[4.5pt]\textbf{C.} #3&\textbf{D.} #4\end{tabularx}}
\newcommand{\one}[4]{\par\noindent\begin{tabularx}{\linewidth}{@{}>{\raggedright\arraybackslash}X@{}}\textbf{A.} #1\\[4pt]\textbf{B.} #2\\[4pt]\textbf{C.} #3\\[4pt]\textbf{D.} #4\end{tabularx}}
\newcommand{\answerbox}{\hfill\textit{Đáp số:}\ \fbox{\rule{0pt}{9pt}\hspace{2.2cm}}}
\setlist[enumerate,1]{label=\textbf{Câu \arabic*.},leftmargin=1.1cm,labelsep=.1cm,itemsep=0.2pt,topsep=0.2pt,parsep=0pt}
\setlist[enumerate,2]{label=\alph*),leftmargin=.55cm,itemsep=1.5pt,topsep=0.2pt,parsep=0pt}

\begin{document}
\noindent
\begin{minipage}[c]{0.38\textwidth}
    \centering
    \textbf{LÀM TOÁN ĐỂ HẠNH PHÚC}\\[-2pt]
    \rule{3.2cm}{0.6pt}\\[3pt]
    \textbf{ĐỀ CHÍNH THỨC}
\end{minipage}\hfill
\begin{minipage}[c]{0.35\textwidth}
    \centering
    \textbf{ĐỀ ÔN THI GIỮA KÌ I}\\
    \textbf{NĂM HỌC 2026 -- 2027}\\[2pt]
    \textbf{ĐỀ SỐ 31}\\
    \textbf{Môn: TOÁN -- Khối lớp 10}\\
    \textit{\small Thời gian làm bài: 90 phút (không kể thời gian phát đề)}
\end{minipage}

\vspace{0.03cm}
\noindent
\begin{minipage}[b]{0.78\textwidth}
    \textbf{Họ và tên:}\dotfill\textbf{SBD:}\makebox[3.5cm]{\dotfill}
\end{minipage}\hfill
\begin{minipage}[b]{0.20\textwidth}
    \hfill
    \setlength{\fboxrule}{0.8pt}%
    \fbox{\begin{minipage}{2.3cm}
        \centering
        \textbf{\small Mã đề thi}\\
        \textbf{0031}
    \end{minipage}}
\end{minipage}
\vspace{0.03cm}

\parttitle{PHẦN I. Trắc nghiệm nhiều lựa chọn}
\textit{Thí sinh trả lời từ câu 1 đến câu 12. Mỗi câu hỏi thí sinh chỉ chọn một phương án.}
\begin{enumerate}
\item Trong các câu sau, câu nào là một mệnh đề toán học?
\two{Tổng ba góc trong một tam giác bằng $180^\circ$.}{Hôm nay trời đẹp quá!}{Bạn có thích học môn Toán không?}{Hãy làm bài tập về nhà!}

\item Cho hai mệnh đề $P$ và $Q$. Phát biểu nào sau đây là đúng về mệnh đề kéo theo $P \Rightarrow Q$?
\two{Mệnh đề $P \Rightarrow Q$ sai khi $P$ sai và $Q$ đúng.}{$P$ là điều kiện đủ để có $Q$.}{$P$ là điều kiện cần để có $Q$.}{Mệnh đề $P \Rightarrow Q$ sai khi cả $P$ và $Q$ cùng sai.}

\item Cho hai tập hợp $A = \{1; 2; 3; 4; 5\}$ và $B = \{3; 4; 5; 6; 7\}$. Khẳng định nào sau đây đúng?
\four{$A \cap B = \{1; 2\}$.}{$A \cup B = \{3; 4; 5\}$.}{$A \setminus B = \{1; 2\}$.}{$B \setminus A = \{1; 2\}$.}

\item Cho hai tập hợp $A = (-\infty; m]$ và $B = [2m - 3; +\infty)$. Có bao nhiêu số nguyên dương $m \in [1; 10]$ để $A \cap B = \varnothing$?
\four{$5$.}{$6$.}{$8$.}{$7$.}

\item Hệ bất phương trình nào sau đây là hệ bất phương trình bậc nhất hai ẩn?
\two{$\begin{cases} 2x - 3y + 1 > 0 \\ x + 4y \le 5 \end{cases}$}{$\begin{cases} x^2 + y \le 4 \\ 2x - y > 1 \end{cases}$}{$\begin{cases} x - 2y + z > 0 \\ y + 3 \le 0 \end{cases}$}{$\begin{cases} xy + 2 > 0 \\ x - y \le 3 \end{cases}$}

\item Điểm nào sau đây thuộc miền nghiệm của bất phương trình bậc nhất hai ẩn $2x - 3y + 6 > 0$?
\four{$M(0; 3)$.}{$N(1; 1)$.}{$P(-4; 0)$.}{$Q(2; 5)$.}

\item Nửa mặt phẳng không bị gạch chéo (kể cả bờ $d$) đi qua hai điểm $(4; 0)$ và $(0; 2)$ chứa gốc toạ độ $O(0; 0)$ là miền nghiệm của bất phương trình nào sau đây?
\four{$x + 2y > 4$.}{$2x + y \le 4$.}{$x + 2y \le 4$.}{$x + 2y \ge 4$.}

\item Diện tích của hình phẳng giới hạn bởi hai trục toạ độ $Ox, Oy$ và miền nghiệm của bất phương trình $3x + 4y \le 12$ ($x \ge 0, y \ge 0$) bằng bao nhiêu?
\four{$12$.}{$8$.}{$7$.}{$6$.}

\item Cho tam giác $ABC$ có các cạnh $BC = a, CA = b, AB = c$. Khẳng định nào sau đây đúng?
\two{$a^2 = b^2 + c^2 - 2bc \cos A$.}{$a^2 = b^2 + c^2 + 2bc \cos A$.}{$a^2 = b^2 + c^2 - 2bc \sin A$.}{$a^2 = b^2 + c^2 + 2bc \sin A$.}

\item Cho tam giác $ABC$ có các cạnh $a, b, c$ và bán kính đường tròn ngoại tiếp $R$. Công thức nào sau đây dùng để tính diện tích $S$ của tam giác $ABC$?
\four{$S = \dfrac{abc}{2R}$.}{$S = \dfrac{abc}{4R}$.}{$S = \dfrac{abc}{R}$.}{$S = 4R \cdot abc$.}

\item Cho góc tù $\alpha$ thoả mãn $\cos \alpha = -\dfrac{3}{5}$. Giá trị của biểu thức $P = \dfrac{\sin \alpha + \tan \alpha}{1 + \cos \alpha}$ bằng
\four{$\dfrac{4}{3}$.}{$-\dfrac{2}{3}$.}{$-\dfrac{4}{3}$.}{$\dfrac{2}{3}$.}

\item Tam giác $ABC$ có độ dài hai cạnh $b = 7, c = 5$ và góc $\widehat{A} = 60^\circ$. Bán kính đường tròn ngoại tiếp $R$ của tam giác $ABC$ bằng
\four{$\sqrt{39}$.}{$2\sqrt{13}$.}{$\dfrac{\sqrt{13}}{2}$.}{$\sqrt{13}$.}
\end{enumerate}
\newpage

\parttitle{PHẦN II. Trắc nghiệm đúng/sai}
\textit{Thí sinh trả lời từ câu 1 đến câu 2. Trong mỗi ý a), b), c), d) ở mỗi câu, thí sinh chọn đúng hoặc sai.}
\begin{enumerate}
\item Cho hai tập hợp $A = \{x \in \mathbb{R} \mid x^2 - 5x + 6 = 0\}$, $B = \{x \in \mathbb{Z} \mid -2 \le x < 3\}$ và tập hợp $C = [m; m + 3]$ (với $m$ là tham số thực).
\begin{enumerate}
    \item Số phần tử của tập hợp $A$ bằng $2$.
    \item Tập hợp $B$ có tất cả $32$ tập hợp con.
    \item Tập hợp hiệu $A \setminus B = \{3\}$.
    \item Có đúng $3$ giá trị nguyên của tham số $m$ để $A \cap C \ne \varnothing$.
\end{enumerate}

\item Để xác định vị trí của một trạm phát sóng $A$ và hai tàu đánh cá $B, C$ trên biển, các kiểm ngư viên đo được khoảng cách $AB = 60\text{ m}$, $AC = 100\text{ m}$ và góc nhìn giữa hai tàu từ trạm phát sóng là $\widehat{BAC} = 120^\circ$.
\begin{enumerate}
    \item Trong tam giác $ABC$, giá trị $\cos \widehat{BAC} = -\dfrac{1}{2}$.
    \item Khoảng cách giữa hai tàu đánh cá $B$ và $C$ bằng $150\text{ m}$.
    \item Diện tích tam giác $ABC$ bằng $1500\sqrt{3}\text{ m}^2$.
    \item Một ca nô tuần tra đi từ trạm $A$ đến một vị trí $M$ trên đoạn thẳng $BC$ sao cho khoảng cách $AM$ ngắn nhất. Khi đó độ dài đoạn thẳng $AM = \dfrac{150\sqrt{3}}{7}\text{ m}$.
\end{enumerate}
\end{enumerate}

\parttitle{PHẦN III. Trắc nghiệm trả lời ngắn}
\textit{Thí sinh trả lời từ câu 1 đến câu 4. Điền kết quả vào ô tương ứng.}
\begin{enumerate}
\item Cho miền tam giác $OAB$ (kể cả biên) với toạ độ các đỉnh là gốc toạ độ $O(0; 0)$, $A(4; 0)$ và $B(1; 3)$ như hình vẽ.

\begin{center}
\vspace{-4pt}
\begin{tikzpicture}[scale=0.75, >=Stealth]
    \draw[line width=1.1pt, ->] (-0.8,0) -- (5.2,0) node[right, font=\small\bfseries] {$x$};
    \draw[line width=1.1pt, ->] (0,-0.8) -- (0,4.2) node[above, font=\small\bfseries] {$y$};
    \node[below left=2pt, font=\small\bfseries] at (0,0) {$O$};
    
    \fill[pattern=north east lines, pattern color=blue!50] (0,0) -- (4,0) -- (1,3) -- cycle;
    \draw[line width=1.3pt, ink] (0,0) -- (4,0) -- (1,3) -- cycle;
    
    \filldraw[black] (4,0) circle (2.4pt) node[below=2pt, font=\small\bfseries] {$A(4;0)$};
    \filldraw[black] (1,3) circle (2.4pt) node[above=2pt, font=\small\bfseries] {$B(1;3)$};
    \draw[dashed, thin] (1,0) -- (1,3) -- (0,3);
    \node[below=2pt, font=\small] at (1,0) {$1$};
    \node[left=2pt, font=\small] at (0,3) {$3$};
\end{tikzpicture}
\vspace{-6pt}
\end{center}

Tìm giá trị nhỏ nhất của tham số $m$ để bất phương trình $2x + y \le m$ nghiệm đúng với mọi điểm $(x; y)$ thuộc miền tam giác $OAB$.
\answerbox

\item Cho tam giác $ABC$ có độ dài ba cạnh lần lượt là $a = 13$, $b = 14$, $c = 15$. Tính bán kính đường tròn nội tiếp $r$ của tam giác $ABC$.
\answerbox
\end{enumerate}
\newpage

\begin{enumerate}[resume]
\item Một xưởng cơ khí sản xuất hai loại chi tiết máy $X$ và $Y$. Mỗi chi tiết $X$ cần $2$ giờ tiện và $1$ giờ phay, đem lại tiền lãi $40$ nghìn đồng. Mỗi chi tiết $Y$ cần $1$ giờ tiện và $3$ giờ phay, đem lại tiền lãi $50$ nghìn đồng. Thời gian tối đa của máy tiện là $180$ giờ/tháng, máy phay là $240$ giờ/tháng. Ngoài ra do nhu cầu thị trường, số lượng chi tiết $Y$ sản xuất không vượt quá $2$ lần số chi tiết $X$. Tìm số tiền lãi lớn nhất mà xưởng có thể đạt được trong một tháng (đơn vị tính: triệu đồng, viết số thập phân).
\answerbox

\item Một trường THPT có $100$ học sinh khối 10 tham gia câu lạc bộ năng khiếu: có $45$ học sinh tham gia CLB Bóng rổ, $53$ học sinh tham gia CLB Cầu lông, và $18$ học sinh không tham gia cả hai câu lạc bộ trên. Hỏi có bao nhiêu học sinh tham gia cả hai CLB Bóng rổ và Cầu lông?
\answerbox
\end{enumerate}

\parttitle{PHẦN IV. Tự luận}
\textit{Thí sinh trình bày chi tiết lời giải từ câu 1 đến câu 4.}
\begin{enumerate}
\item Để đo khoảng cách giữa hai vị trí $A$ và $B$ bị ngăn cách bởi một hồ nước sâu mà không thể đo trực tiếp, người ta chọn một điểm mốc $C$ quan sát được cả $A$ và $B$. Bằng giác kế, người ta đo được khoảng cách $CA = 120\text{ m}$, $CB = 150\text{ m}$ và góc nhìn $\widehat{ACB} = 60^\circ$.
\begin{enumerate}
    \item Tính khoảng cách $AB$ giữa hai vị trí (làm tròn kết quả đến hàng phần mười theo đơn vị mét).
    \item Tính diện tích tam giác $ABC$.
\end{enumerate}

\item Cho hai tập hợp $A = (-4; 5]$ và $B = [m - 1; 2m + 3)$ (với $m$ là tham số thực).
\begin{enumerate}
    \item Khi $m = 2$, hãy xác định các tập hợp $A \cap B$, $A \cup B$ và $A \setminus B$.
    \item Tìm tất cả các giá trị thực của tham số $m$ để tập hợp $B$ là tập con khác rỗng của tập hợp $A$.
\end{enumerate}

\item Một trang trại dự định trồng hai loại cây ăn quả: Cây cam và Cây bưởi trên tổng diện tích đất tối đa là $12\text{ ha}$.
\begin{itemize}[leftmargin=1.5em, itemsep=0.1pt]
    \item Để trồng $1\text{ ha}$ cam cần $20$ triệu đồng vốn đầu tư và $30$ công lao động.
    \item Để trồng $1\text{ ha}$ bưởi cần $30$ triệu đồng vốn đầu tư và $20$ công lao động.
\end{itemize}
Trang trại hiện có số vốn đầu tư tối đa là $270$ triệu đồng và tổng số công lao động tối đa có thể huy động là $300$ công. Sau khi thu hoạch, mỗi héc-ta cam mang lại lợi nhuận $40$ triệu đồng, mỗi héc-ta bưởi mang lại lợi nhuận $50$ triệu đồng.
\begin{enumerate}
    \item Gọi $x, y$ (héc-ta) lần lượt là diện tích trồng cam và bưởi ($x \ge 0, y \ge 0$). Hãy lập hệ bất phương trình bậc nhất hai ẩn thể hiện các điều kiện ràng buộc trên và biểu diễn miền nghiệm của hệ trên mặt phẳng toạ độ $Oxy$.
    \item Xác định diện tích trồng mỗi loại cây để trang trại thu được lợi nhuận lớn nhất. Số tiền lợi nhuận lớn nhất đó bằng bao nhiêu?
\end{enumerate}

\item Cho hai tập hợp khác rỗng $A = [m - 1; m + 2)$ và $B = (-\infty; -2) \cup [4; +\infty)$ với $m$ là tham số thực.
\begin{enumerate}
    \item Tìm điều kiện của tham số $m$ để $A \cap B = \varnothing$.
    \item Tìm tất cả các giá trị nguyên của $m$ để tập hợp $C = A \setminus B$ có đúng $3$ phần tử là số nguyên.
\end{enumerate}
\end{enumerate}

\newpage

% ==================== ĐÁP ÁN VÀ LỜI GIẢI CHI TIẾT ====================
\parttitle{BẢNG ĐÁP ÁN TỔNG HỢP}

\vspace{0.2cm}
\noindent\textbf{PHẦN I. Trắc nghiệm nhiều lựa chọn (12 câu)}
\vspace{0.15cm}

\begin{center}
\begin{tabular}{|c|c|c|c|c|c|c|c|c|c|c|c|}
\hline
\textbf{1} & \textbf{2} & \textbf{3} & \textbf{4} & \textbf{5} & \textbf{6} & \textbf{7} & \textbf{8} & \textbf{9} & \textbf{10} & \textbf{11} & \textbf{12} \\
\hline
\textbf{A} & \textbf{B} & \textbf{C} & \textbf{D} & \textbf{A} & \textbf{B} & \textbf{C} & \textbf{D} & \textbf{A} & \textbf{B} & \textbf{C} & \textbf{D} \\
\hline
\end{tabular}
\end{center}

\vspace{0.3cm}
\noindent\textbf{PHẦN II. Trắc nghiệm đúng/sai (2 câu)}
\vspace{0.15cm}

\begin{center}
\begin{tabular}{|c|c|c|c|c|}
\hline
\textbf{Câu} & \textbf{Lệnh hỏi a)} & \textbf{Lệnh hỏi b)} & \textbf{Lệnh hỏi c)} & \textbf{Lệnh hỏi d)} \\
\hline
\textbf{Câu 1} & \textbf{Đúng} & \textbf{Đúng} & \textbf{Đúng} & \textbf{Sai} \\
\hline
\textbf{Câu 2} & \textbf{Đúng} & \textbf{Sai} & \textbf{Đúng} & \textbf{Đúng} \\
\hline
\end{tabular}
\end{center}

\vspace{0.3cm}
\noindent\textbf{PHẦN III. Trắc nghiệm trả lời ngắn (4 câu)}
\vspace{0.15cm}

\begin{center}
\begin{tabular}{|c|c|c|c|c|}
\hline
\textbf{Câu} & \textbf{Câu 1} & \textbf{Câu 2} & \textbf{Câu 3} & \textbf{Câu 4} \\
\hline
\textbf{Đáp số} & $8$ & $4$ & $5{,}4$ & $16$ \\
\hline
\end{tabular}
\end{center}

\vspace{0.4cm}
\parttitle{PHẦN IV. ĐÁP SỐ TỰ LUẬN}
\vspace{0.15cm}

\begin{enumerate}[leftmargin=1.5em, itemsep=6pt]
    \item[\textbf{Câu 1.}] 
    \begin{itemize}[leftmargin=1.2em, itemsep=2pt]
        \item[a)] $AB = 30\sqrt{21} \approx 137{,}5\text{ m}$.
        \item[b)] $S_{ABC} = 4500\sqrt{3} \approx 7794{,}2\text{ m}^2$.
    \end{itemize}

    \item[\textbf{Câu 2.}] 
    \begin{itemize}[leftmargin=1.2em, itemsep=2pt]
        \item[a)] $A \cap B = [1; 5]$, \quad $A \cup B = (-4; 7)$, \quad $A \setminus B = (-4; 1)$.
        \item[b)] $-3 < m \le 1$.
    \end{itemize}

    \item[\textbf{Câu 3.}] 
    \begin{itemize}[leftmargin=1.2em, itemsep=2pt]
        \item[a)] Hệ bất phương trình mô hình hóa diện tích canh tác, vốn và công lao động:
        \[
        \begin{cases}
        x \ge 0,\ y \ge 0 \\
        x + y \le 12 \\
        2x + 3y \le 27 \\
        3x + 2y \le 30
        \end{cases}
        \]
        \item[b)] Trang trại nên trồng $7{,}2\text{ ha}$ cam và $4{,}2\text{ ha}$ bưởi để thu được lợi nhuận lớn nhất là $498$ triệu đồng.
    \end{itemize}

    \item[\textbf{Câu 4.}] 
    \begin{itemize}[leftmargin=1.2em, itemsep=2pt]
        \item[a)] $-1 \le m \le 2$.
        \item[b)] Các giá trị nguyên của $m$ là $m \in \{-1; 0; 1; 2\}$ (có đúng $4$ giá trị nguyên).
    \end{itemize}
\end{enumerate}

\newpage

% ==================== HÌNH VẼ MIỀN NGHIỆM HỆ BẤT PHƯƠNG TRÌNH ====================
\parttitle{HÌNH VẼ MIỀN NGHIỆM HỆ BẤT PHƯƠNG TRÌNH}

\vspace{0.2cm}
\noindent\textbf{1. Miền nghiệm Câu 3 -- Phần III (Bài toán xưởng cơ khí sản xuất hai chi tiết máy):}
\[
\text{Hệ bất phương trình ràng buộc: }
\begin{cases}
2x + y \le 180 \\
x + 3y \le 240 \\
y \le 2x \\
x \ge 0,\ y \ge 0
\end{cases}
\]

\begin{center}
\begin{tikzpicture}[x=0.092cm, y=0.092cm]
    % Lưới toạ độ phụ
    \draw[help lines, color=gray!25, step=20] (-10,-10) grid (115,95);
    
    % Trục tọa độ
    \draw[-{Latex[length=3mm]}, line width=1.1pt] (-12,0) -- (120,0) node[below=2pt, font=\normalsize\bfseries] {$x$};
    \draw[-{Latex[length=3mm]}, line width=1.1pt] (0,-12) -- (0,105) node[left=2pt, font=\normalsize\bfseries] {$y$};
    \node[below left=2pt, font=\normalsize\bfseries] at (0,0) {$O$};
    
    % Tô màu miền nghiệm OABC
    \fill[ink!20] (0,0) -- (90,0) -- (60,60) -- (34.286,68.571) -- cycle;
    \draw[line width=1.4pt, ink] (0,0) -- (90,0) -- (60,60) -- (34.286,68.571) -- cycle;
    
    % Các đường thẳng biên
    % d1: 2x + y = 180
    \draw[line width=1.1pt, blue!80!black] (42,96) -- (96,-12) 
        node[pos=0.02, above right, font=\small\bfseries] {$d_1: 2x + y = 180$};
        
    % d2: x + 3y = 240
    \draw[line width=1.1pt, red!80!black] (-10,83.33) -- (118,40.67) 
        node[pos=0.95, above right, font=\small\bfseries] {$d_2: x + 3y = 240$};
        
    % d3: y = 2x
    \draw[line width=1.1pt, teal!80!black] (-5,-10) -- (48,96) 
        node[pos=0.98, above left, font=\small\bfseries] {$d_3: y = 2x$};

    % Các đường gióng toạ độ
    \draw[dashed, line width=0.8pt, gray!80] (60,0) node[below=2pt, font=\small\bfseries] {$60$} -- (60,60) -- (0,60) node[left=2pt, font=\small\bfseries] {$60$};
    \draw[dashed, line width=0.8pt, gray!80] (34.286,0) node[below=2pt, font=\small\bfseries] {$\frac{240}{7}$} -- (34.286,68.571) -- (0,68.571) node[left=2pt, font=\small\bfseries] {$\frac{480}{7}$};
    \node[below=2pt, font=\small\bfseries] at (90,0) {$90$};
    \node[left=2pt, font=\small\bfseries] at (0,80) {$80$};

    % Các điểm đỉnh
    \fill[black] (0,0) circle (2.5pt);
    \fill[black] (90,0) circle (2.5pt) node[above right=1pt, font=\normalsize\bfseries] {$A(90;0)$};
    \fill[black] (60,60) circle (2.5pt) node[above right=1pt, font=\normalsize\bfseries] {$B(60;60)$};
    \fill[black] (34.286,68.571) circle (2.5pt) node[above left=1pt, font=\normalsize\bfseries] {$C\left(\frac{240}{7};\frac{480}{7}\right)$};

    % Nhãn miền nghiệm
    \node[font=\bfseries\color{ink}] at (48,28) {Miền nghiệm $OABC$};
\end{tikzpicture}
\end{center}

\vspace{0.2cm}
\begin{itemize}[leftmargin=1.5em, itemsep=3pt]
    \item Miền nghiệm là tứ giác $OABC$ (kể cả biên) với toạ độ các đỉnh:
    \[
    O(0;0), \quad A(90;0), \quad B(60;60), \quad C\left(\frac{240}{7}; \frac{480}{7}\right).
    \]
    \item Hàm tiền lãi: $L(x, y) = 40x + 50y$ (nghìn đồng).
    \item Giá trị tại các đỉnh: $L(O) = 0$; $L(A) = 3600$; $L(B) = 5400$; $L(C) \approx 4800$.
    \item \textbf{Kết luận:} Tiền lãi lớn nhất là $5400$ nghìn đồng $= \mathbf{5{,}4}$ triệu đồng (đạt tại đỉnh $B(60;60)$).
\end{itemize}

\newpage

% ==================== MIỀN NGHIỆM CÂU 3 PHẦN IV ====================
\parttitle{HÌNH VẼ MIỀN NGHIỆM HỆ BẤT PHƯƠNG TRÌNH (Tiếp theo)}

\vspace{0.2cm}
\noindent\textbf{2. Miền nghiệm Câu 3 -- Phần IV (Bài toán quy hoạch trang trại trồng cam và bưởi):}
\[
\text{Hệ bất phương trình ràng buộc: }
\begin{cases}
x \ge 0,\ y \ge 0 \\
x + y \le 12 \\
2x + 3y \le 27 \\
3x + 2y \le 30
\end{cases}
\]

\begin{center}
\begin{tikzpicture}[x=0.82cm, y=0.82cm]
    % Lưới toạ độ
    \draw[help lines, color=gray!25, step=2] (-1,-1) grid (14,13);
    
    % Trục tọa độ
    \draw[-{Latex[length=3mm]}, line width=1.1pt] (-1.2,0) -- (14.5,0) node[below=2pt, font=\normalsize\bfseries] {$x$};
    \draw[-{Latex[length=3mm]}, line width=1.1pt] (0,-1.2) -- (0,13.5) node[left=2pt, font=\normalsize\bfseries] {$y$};
    \node[below left=2pt, font=\normalsize\bfseries] at (0,0) {$O$};

    % Tô màu miền nghiệm OABC
    \fill[ink!20] (0,0) -- (10,0) -- (7.2,4.2) -- (0,9) -- cycle;
    \draw[line width=1.4pt, ink] (0,0) -- (10,0) -- (7.2,4.2) -- (0,9) -- cycle;

    % d1: 3x + 2y = 30
    \draw[line width=1.1pt, blue!80!black] (2,12) -- (10.67,-1) 
        node[pos=0.08, above right, font=\small\bfseries] {$d_1: 3x + 2y = 30$};

    % d2: 2x + 3y = 27
    \draw[line width=1.1pt, red!80!black] (-1,9.67) -- (14,-0.33) 
        node[pos=0.92, above right, font=\small\bfseries] {$d_2: 2x + 3y = 27$};

    % d3: x + y = 12 (đường biên nằm ngoài miền nghiệm)
    \draw[line width=1.1pt, dashed, orange!90!black] (-0.5,12.5) -- (13,-1) 
        node[pos=0.25, above right, font=\small\bfseries] {$d_3: x + y = 12$};

    % Đường gióng toạ độ
    \draw[dashed, line width=0.8pt, gray!80] (7.2,0) node[below=2pt, font=\small\bfseries] {$7{,}2$} -- (7.2,4.2) -- (0,4.2) node[left=2pt, font=\small\bfseries] {$4{,}2$};
    \node[below=2pt, font=\small\bfseries] at (10,0) {$10$};
    \node[left=2pt, font=\small\bfseries] at (0,9) {$9$};

    % Các đỉnh
    \fill[black] (0,0) circle (2.5pt);
    \fill[black] (10,0) circle (2.5pt) node[above right=1pt, font=\normalsize\bfseries] {$A(10;0)$};
    \fill[black] (7.2,4.2) circle (2.5pt) node[above right=1pt, font=\normalsize\bfseries] {$B(7{,}2; 4{,}2)$};
    \fill[black] (0,9) circle (2.5pt) node[above right=1pt, font=\normalsize\bfseries] {$C(0;9)$};

    % Nhãn miền nghiệm
    \node[font=\bfseries\color{ink}] at (3.8,3.2) {Miền nghiệm $OABC$};
\end{tikzpicture}
\end{center}

\vspace{0.2cm}
\begin{itemize}[leftmargin=1.5em, itemsep=3pt]
    \item Miền nghiệm là tứ giác $OABC$ (kể cả biên) với toạ độ các đỉnh:
    \[
    O(0;0), \quad A(10;0), \quad B(7{,}2; 4{,}2), \quad C(0;9).
    \]
    *(Đường thẳng $x + y = 12$ nằm ngoài tứ giác $OABC$ nên là ràng buộc dư thừa).*
    \item Hàm lợi nhuận: $F(x, y) = 40x + 50y$ (triệu đồng).
    \item Giá trị tại các đỉnh: $F(O) = 0$; $F(A) = 400$; $F(B) = 498$; $F(C) = 450$.
    \item \textbf{Kết luận:} Lợi nhuận lớn nhất là $\mathbf{498}$ triệu đồng khi trồng $\mathbf{7{,}2\text{ ha}}$ cam và $\mathbf{4{,}2\text{ ha}}$ bưởi.
\end{itemize}

\end{document}
